\chapter{Discussion}\label{chap:discussion}

\subsection{Summary of Contributions}
Recap the main results in 1--2 pages.

\subsection{Limitations}

\subsubsection{Modeling Limitations}
- Fixed truth labels and unique dependency structures (real proofs have multiple strategies).
- Finite formula universe (real mathematics is open-ended).
- Simplified proof/conjecture kernels (real LLM provers have complex failure modes).
- No modeling of communication beyond publication (real agents could discuss strategies).

\subsubsection{Theoretical Limitations}
- Convergence guarantees are weak (meta-game CCE, not full-game Nash or even CCE).
- The gap between meta-game equilibria and full-game equilibria is not characterized.
- PPAD and NEXP hardness are worst-case; average-case complexity is unknown.

\subsubsection{Empirical Limitations}
- Synthetic experiments on relatively small instances; scalability to real-world size is not demonstrated.
- The Lean experiments use a limited set of LLM provers; diversity is artificially constrained.
- Market mechanisms are tested in isolation; hybrid mechanisms are not explored.

\subsection{Future Work}

\subsubsection{Theoretical Extensions}
- Sharper characterization of the cost of decentralization (price-of-anarchy-style bounds for specific VAMP classes).
- Average-case complexity analysis for natural distributions over VAMP instances.
- Mechanism design: optimal market mechanism selection given a VAMP instance class.

\subsubsection{Architectural Extensions}
- Graph neural network observation encoders for better scalability.
- Hierarchical policies for long-horizon credit assignment.
- Communication channels between agents (beyond publication).

\subsubsection{Applications Beyond Theorem Proving}
- Formal verification of software (proof obligations = verification conditions).
- Drug discovery (tasks = biological assays, agents = heterogeneous screening platforms).
- Collaborative AI research (agents = different LLMs with different training, tasks = research sub-questions).

\subsection{Broader Impact}
Discuss the implications of market-mediated AI coordination: potential for more efficient scientific discovery, risks of market manipulation by sophisticated agents, fairness considerations, and the relationship between market efficiency and social welfare.
